\begin{abstract}
Cylinder chess joins the a-file to the h-file, so the board has no edges or corners. We study its
king-and-pawn endgames with combinatorial game theory, following Elkies: a passed pawn counts as a win, and a
mutual zugzwang of the kings counts as zero. We first repair this model. We replace the usual passed-pawn test,
which ignores en passant, with one that respects it, and we give conditions under which frozen kings are
exact. Under the new test, an exhaustive census of small positions finds no nimber above star with six pawns or
fewer, but star two appears with seven pawns, in two positions confirmed by three independent programs. We also
find one-eighth, up-second star and Tiny 1, and test the model against real chess with
endgame tablebases. For the one-eighth position, a queen-promotion search leaves one undecided end
position, with black to move, on which its value depends. Four-man tablebases on the cylinder and the torus show
that king and rook can no longer force mate. We correct the earlier analysis of Dawson's chess on the cylinder
and prove that the full ring is a first-player win exactly when the strip three files shorter has Grundy value
zero.
\end{abstract}

\begin{IEEEkeywords}
Chess, combinatorial game theory, endgame tablebases, game solving, nimbers.
\end{IEEEkeywords}
